\section*{Chapitre \ref{ch:s2:supply-trades} --- Trades d'offre}

\begin{solution}{exo:s2:supply-trades:1}
$0.07 \times 40 = 2.8$ points de base.
\end{solution}

\begin{solution}{exo:s2:supply-trades:2}
$3.44/17 \times \sqrt{12} = 0.70$.
\end{solution}

\begin{solution}{exo:s2:supply-trades:3}
$0.0009 \times \$40$ milliards $= \$36$ millions à $0.0018 \times \$40$ milliards $= \$72$ millions.
\end{solution}

\begin{solution}{exo:s2:supply-trades:4}
Connaître l'offre ne revient pas à disposer du bilan nécessaire pour la porter~: les dealers et les investisseurs doivent faire de la place aux nouvelles obligations, et ceux qui
disposent de capacités la font payer. Les investisseurs finaux arrivent lentement, donc le prix bouge jusqu'à ce qu'ils soient là.
\end{solution}

\begin{solution}{exo:s2:supply-trades:5}
Le bilan des dealers est devenu plus coûteux depuis la crise~: ils prennent des positions plus petites et les replacent plus vite~; les fonds d'investissement enchérissent désormais davantage lors des
adjudications, en se partageant l'offre que les dealers absorbaient autrefois.
\end{solution}

\begin{solution}{exo:s2:supply-trades:6}
L'adjudication se clôt à 13~h, et le rendement de fin de journée intègre déjà une partie de la reprise qui la suit~; le
relèvement de rendement précédant l'adjudication est en partie masqué.
\end{solution}

\begin{solution}{exo:s2:supply-trades:7}
La jambe d'avant reste inchangée (1.71 point de base par adjudication)~; la jambe d'après ne rapporte que ses coûts ($-0.09$)~; les deux ensemble 1.62. La \omterm{def:s2:supply-trades:concession}{concession d'adjudication}
payée avant l'adjudication est ce qui rémunère les dealers~; la reprise qui suit est un pari distinct sur l'achat des investisseurs finaux.
\end{solution}

\begin{solution}{exo:s2:supply-trades:8}
Prises isolément, les adjudications ont un écart-type d'environ 17 points de base pour une moyenne de 3.44~; trois pertes de 20 restent largement dans le bruit. Il faut
des dizaines d'adjudications pour détecter l'effet, et autant pour détecter sa disparition.
\end{solution}

\begin{solution}{pb:s2:supply-trades:1}
\begin{enumerate}
\item Le relèvement de rendement avant une nouvelle émission, repris après~; la \omterm{def:s2:supply-trades:otr}{prime on-the-run} de la dernière émission à une échéance donnée.
\item Les cours baissent avant les adjudications et remontent après~; un coût caché de 9 à 18 points de base selon la taille.
\item Les émissions gonflent les inventaires des dealers~; ceux-ci sont rémunérés par la hausse ultérieure des cours~; cette rémunération a baissé avec la hausse des \omterm{def:s2:swap-spread-and-asset-swap-trades:bscost}{coûts de bilan} et la participation accrue des
fonds.
\item Ils s'échangent avec une prime au titre de la liquidité et de la valeur en repo.
\item Regroupées, 1.47 points de base de hausse le jour de l'adjudication par rapport à dix jours avant, contre 0.40 cinq jours après.
\item Les titres à 10 et 30 ans voient leur rendement monter avant~; ceux à 5 et 7 ans le voient baisser après~; les 2 ans ne montrent rien.
\item La hausse d'avant croît avec la taille (1.48, 1.69, 2.95 points de base par tercile)~; la baisse d'après non.
\item La partie de la journée d'adjudication précédant 13~h.
\item Position courte sur l'offre connue avant la vente, position longue après.
\item 0.07 point de base par milliard à capacité moyenne, 70~\% restitués après.
\item 3.44 par adjudication, $t$ = 3.1, ratio de Sharpe 0.70~; chaque jambe environ 1.7 et 0.50.
\item Un bruit quotidien de 5.5 points de base sur dix jours noie une \omterm{def:s2:supply-trades:concession}{concession d'adjudication} de 2.6.
\item \textbf{Résultat nommé.} Le trade d'adjudication rapporte 3.44 points de base de rendement par adjudication ($t$ = 3.1, ratio de Sharpe 0.70 par an), et chaque milliard de
taille ajoute 0.24 point de base en moyenne, bien que le résultat d'une adjudication isolée reste du bruit.
\item Autant que la taille~: la concession est la taille rapportée à la capacité.
\item Autour du calendrier des grandes émissions, couvertes par des Treasuries ou des swaps, en prenant les allocations à la \omterm{def:s2:supply-trades:concession}{concession d'adjudication}.
\item Acheter de la duration avant l'allongement de fin de mois de l'indice, pour la demande des trackers d'indice.
\item Rendements aux heures de trade, coûts par jambe, de nombreuses adjudications.
\item Le roll on-the-run.
\item La courbe ajustée du chapitre~10 exclut les émissions autour desquelles ce chapitre fait des trades.
\item L'émetteur, par la \omterm{def:s2:supply-trades:concession}{concession d'adjudication} qu'il paie pour un placement rapide.
\end{enumerate}
\end{solution}

\begin{solution}{iq:s2:supply-trades:1}
Les dealers doivent absorber la nouvelle offre avec un bilan limité, et les investisseurs finaux arrivent lentement~; le marché intègre dans le prix le coût de détention de
l'offre jusqu'à sa distribution.
\end{solution}

\begin{solution}{iq:s2:supply-trades:2}
Mesurer les variations de rendement du titre adjugé et de ses substituts proches sur des fenêtres précédant et suivant chaque adjudication, par rapport à un indice de référence ou à une
courbe ajustée~; contrôler les actualités macroéconomiques, la fin de mois, les autres adjudications voisines et le niveau des taux.
\end{solution}

\begin{solution}{iq:s2:supply-trades:3}
Prendre une position courte sur des contrats à terme ou sur le titre when-issued avant l'adjudication, en proportion de l'attribution attendue, et ajuster après le résultat~; prévoir la
distribution des obligations aux clients.
\end{solution}

\begin{solution}{iq:s2:supply-trades:4}
Les actualités macroéconomiques autour des dates d'adjudication dominent chaque adjudication~; l'effet peut s'atténuer à mesure que davantage de desks le trade~; les coûts et l'encombrement (crowding) sur les contrats à terme à l'approche de
l'adjudication.
\end{solution}

\begin{solution}{iq:s2:supply-trades:5}
Annonces (date, taille, titre), résultats (high yield, bid-to-cover, allocations) et cotations when-issued, horodatées~; contrôles sur les réouvertures
et les changements du calendrier~; une base indexée par CUSIP et par date d'adjudication.
\end{solution}

\begin{solution}{iq:s2:supply-trades:6}
Sur une fenêtre de $k$ jours, le bruit vaut $\sigma\sqrt k$~; avec $n$ adjudications, l'erreur type de la moyenne vaut $\sigma\sqrt{k/n}$, donc $t = c\sqrt n/(\sigma\sqrt
k) = 2$ donne $n = 4k\sigma^2/c^2$~: avec $\sigma = 5.5$, $k = 5$ et $c = 2.6$, environ 90 adjudications pour la jambe d'avant.
\end{solution}
